% Solutions of Book 4, chapter 11 (processor-and-pipelining).
\section*{Chapter \ref{ch:b2:processor-and-pipelining} --- The Processor: Datapath, Control and Pipelining}

\begin{solution}{exo:b2:processor-and-pipelining:1}
By \cref{prop:b2:processor-and-pipelining:speedup}, $n + 4$ cycles: 12 for 8 instructions. The pipeline has
$12 \times 5 = 60$ stage slots in those cycles and the instructions use $8 \times 5 = 40$: two thirds. The empty
slots are the filling at the start and the draining at the end, which matter less as $n$ grows.
\end{solution}

\begin{solution}{exo:b2:processor-and-pipelining:2}
In an S-type word, bits 11 to 7 are the low five bits of the immediate, not a register. For \code{sw t2, 12(sp)}
they hold 12, so the \omterm{def:b2:assembly-and-calling-convention:register}{register file} writes register \code{x12}, which is \code{a2}. MemtoReg is 0, so the value
written is the ALU result, the address $\code{sp} + 12$. The store itself is correct; \code{a2} is silently
overwritten, and the bug shows only when a caller later reads \code{a2}. That is why RegWrite must be 0 for every
instruction without a destination, while the ``don't care'' signals may be anything.
\end{solution}

\begin{solution}{exo:b2:processor-and-pipelining:3}
\code{lui rd, imm} writes $\mathit{imm} \cdot 2^{12}$ to \code{rd}: RegWrite 1, ALUSrc 1, MemRead 0, MemWrite 0,
MemtoReg 0, Branch 0. Two changes are needed. The immediate generator learns the U format (the upper 20 bits of the
word, low bits zero). And the ALU's first input must be 0, not \code{rs1}, since bits 19 to 15 of a U-type word
belong to the immediate: a new multiplexer on that input, driven by a new signal, with the ALU adding.
\omcode{code/bachelor-2/11-processor-and-pipelining/exercises/datapath_more.py}{4}{10}{The new signals for
\code{lui}, \code{jal} and \code{jalr}}
The chapter's test runs a program with \code{lui} (through \code{li}) and checks the model against \code{rv}.
\end{solution}

\begin{solution}{exo:b2:processor-and-pipelining:4}
For $k = 5$: $50/14 \approx 3.57$ for $n = 10$, $500/104 \approx 4.81$ for $n = 100$, $5000/1004 \approx 4.98$ for
$n = 1000$. For $k = 10$, $n = 10$: $100/19 \approx 5.26$, still far from 10, as filling and draining a deep pipeline
cost more. For 90\,\% of 5: $5n/(n + 4) \ge 4.5$ gives $0.5n \ge 18$, so $n \ge 36$.
\end{solution}

\begin{solution}{exo:b2:processor-and-pipelining:5}
\omoutput{code/bachelor-2/11-processor-and-pipelining/outputs/run_exo_diagram.txt}{The model's diagram.}
The \code{addi t1} uses \code{t0} just after it is loaded: one \omterm{def:b2:processor-and-pipelining:stall}{stall}, by
\cref{prop:b2:processor-and-pipelining:load-use}, and the next \code{addi} waits one cycle in IF. \code{add t3}
gets \code{t1} from MEM/WB and \code{t2} from EX/MEM. Five instructions take $5 + 4 + 1 = 10$ cycles.
\end{solution}

\begin{solution}{exo:b2:processor-and-pipelining:6}
\omoutput[9-13]{code/bachelor-2/11-processor-and-pipelining/outputs/run_forwarding.txt}{The faulty unit on
\code{exo_x0.S}.}
\code{addi zero, t0, 1} computes 42 and, since it writes \code{x0}, the \omterm{def:b2:assembly-and-calling-convention:register}{register file} ignores it. But the faulty unit
sees \code{rd == 0 == rs} in the next instruction and forwards 42 from EX/MEM to both inputs of \code{add a0, zero,
zero}, which computes 84 instead of 0. Writes to \code{x0} are rare in compiled code, but a \code{nop} is
\code{addi zero, zero, 0}, and any instruction can be used as a no-op that way.
\end{solution}

\begin{solution}{exo:b2:processor-and-pipelining:7}
\omcode{code/bachelor-2/11-processor-and-pipelining/exercises/btfn.py}{5}{16}{Backward taken, forward not taken}
The predictor decodes the branch's offset; it needs no memory and no warm-up. On the three programs it mispredicts 1,
72 and 158 branches (two-bit: 2, 60 and 161), for CPIs of 1.279, 1.297 and 1.322. It wins on \code{loop_sum.S}, whose
only branch closes a loop, and on \code{collatz.S}, whose backward \code{bge} closes the outer loop; it loses on the
sort, where the forward branch \code{bge t4, t3, keep} is taken 55 times out of 120, which a counter
follows better than a fixed rule. The test checks the case of \code{loop_sum.S}.
\end{solution}

\begin{solution}{exo:b2:processor-and-pipelining:8}
Let a load be in EX at cycle $x$. In the model a consumer may be in EX at $x + 2$ at the earliest (the word leaves
data memory at the end of MEM, cycle $x + 1$, and reaches the ALU input through MEM/WB). Instructions enter EX in
order, at most one per cycle, and without \omterm{def:b2:processor-and-pipelining:hazard}{hazards} each enters EX one cycle after the one before. If the consumer
follows the load directly, it would be in EX at $x + 1$; the rule holds it one cycle, to $x + 2$, and nothing else
delays it by this dependence: exactly one \omterm{def:b2:processor-and-pipelining:stall}{stall}. If it comes $j \ge 2$ instructions later, it cannot be in EX before
$x + j \ge x + 2$, since the instructions between enter EX in the cycles in between: the rule asks nothing more, no
\omterm{def:b2:processor-and-pipelining:stall}{stall}. (The test checks both cases with an ALU instruction, a store and a branch as the consumer.)
\end{solution}

\begin{solution}{exo:b2:processor-and-pipelining:9}
The \code{addi t0, t0, 4} does not depend on the load; moving it between the load and the add puts the use two
places after the load:
\omcode{code/bachelor-2/11-processor-and-pipelining/exercises/loop_sum_sched.S}{9}{12}{The reordered loop}
The \code{bltu} still reads \code{t0} one instruction after the \code{addi}, which \omterm{def:b2:processor-and-pipelining:forwarding}{forwarding} handles. The model gives
202 cycles instead of 234, the 32 load-use \omterm{def:b2:processor-and-pipelining:stall}{stalls} removed, with the same 136 instructions. Compilers do this
reordering, called instruction scheduling, for exactly this reason.
\end{solution}

\begin{solution}{exo:b2:processor-and-pipelining:10}
Per instruction, the branches cost $0.2 \times 0.1 \times 2 = 0.04$ cycle and the load-uses $0.25 \times 0.4 \times 1
= 0.10$ cycle: $\text{CPI} = 1 + 0.04 + 0.10 = 1.14$. Halving the misprediction rate saves $0.02$ (CPI 1.12); halving
the load-uses saves $0.05$ (CPI 1.09). Here the reordering is worth more, because load-uses are more frequent than
mispredictions, even with a smaller penalty.
\end{solution}

\begin{solution}{exo:b2:processor-and-pipelining:11}
Within one execution the branch is taken $m - 1 \ge 1$ times, then not taken. Suppose the counter is at least 2 when
an execution starts. The taken outcomes are predicted taken and push the counter to 3, where it stays; the final
not-taken outcome is predicted taken, wrongly, and brings the counter down to 2. So the execution has exactly one
misprediction and the next one starts at 2. Once the counter has reached 3 in one execution, that execution ends at
2, and by induction every later one starts at 2 and has one misprediction. With one bit, an execution starts after a
not-taken outcome: its first outcome, taken, is mispredicted, the next $m - 2$ are right, and the exit is wrong
again: two per execution. For $m = 1$ the branch is never taken; the counter falls to 0 and never errs again, while
the hypothesis (reaching 3) is never met: the proposition does not apply, and the loop costs nothing once warm.
\end{solution}

\begin{solution}{exo:b2:processor-and-pipelining:12}
\code{jal rd, imm} writes $\mathit{pc} + 4$ to \code{rd} and jumps to $\mathit{pc} + \mathit{imm}$ (J format);
\code{jalr rd, imm(rs1)} writes $\mathit{pc} + 4$ and jumps to $(\code{rs1} + \mathit{imm})$ with bit 0 cleared. The
datapath needs a third input on the write-back multiplexer, $\mathit{pc} + 4$ from the existing adder (signal
\code{pc4_to_reg}); the J format in the immediate generator, whose result goes to the existing branch adder; and a
third input on the PC multiplexer, the ALU result with bit 0 cleared (signal \code{jump}, three values: none,
\code{jal}, \code{jalr}).
\omcode{code/bachelor-2/11-processor-and-pipelining/exercises/datapath_more.py}{25}{39}{The extended cycle}
The test runs a program with a \code{call}, a \code{ret} and an indirect \code{jalr} through a register: the model
agrees with \code{rv} on every instruction but the \code{auipc} of \code{la} and the final \code{ecall}, and on every
program of the chapter it now covers the \code{ret}s too.
\end{solution}

\begin{solution}{pb:b2:processor-and-pipelining:1}
\textbf{1.} Five without a swap (\code{lw}, \code{lw}, \code{bge}, \code{addi}, \code{bltu}); seven with one, the two
stores added.

\textbf{2.} Six: \code{blez}, \code{mv}, \code{slli}, \code{add}, then \code{addi} and \code{j} after the inner loop.

\textbf{3.} The \code{addi} before the loop runs once. The outer loop runs for $i = n - 1$ down to 1, and turn $i$
executes 6 instructions and $i$ inner turns of 5, plus 2 per swap: $\sum_{i=1}^{n-1} (6 + 5i) = 6(n - 1) +
\frac{5}{2}n(n-1)$, and $2S$ for the swaps. The last test \code{blez}, taken, and the \code{ret} add 2. In all
$\frac{5}{2}n(n-1) + 6(n-1) + 2S + 3$; for $n = 1$, the three instructions \code{addi}, \code{blez}, \code{ret}.

\textbf{4.} A swap exchanges two adjacent values $a_j > a_{j+1}$. The relative order of every other pair is unchanged,
so the number of inversions falls by exactly one. At the end the array is sorted, with no inversion; every swap
removed one, so $S$ equals the initial number of inversions.

\textbf{5.} $600 + 90 + 130 + 3 = 823$ for \code{sort}, and $823 + 141 = 964$, the count of
\cref{out:b2:processor-and-pipelining:cpi}.

\textbf{6.} Loads: 2 per inner turn, 240, the 16 of the printing loop and the \code{lw ra}: 257. Stores: 2 per swap,
130, and \code{sw ra}: 131.

\textbf{7.} \code{bge t4, t3} uses \code{t4}, loaded just before, and \code{t3}, two before; \code{bltu t1, t2} uses
\code{t1}, set by the \code{addi} just before. The two stores use the loaded values two and three places after their
loads.

\textbf{8.} Without \omterm{def:b2:processor-and-pipelining:forwarding}{forwarding} the consumer reads in ID, at the earliest in the producer's WB cycle, $x + 2$, so it is
in EX at $x + 3$. Placed $j$ after, it would naturally be in EX at $x + j$: it waits $3 - j$ cycles, 2 for $j = 1$, 1
for $j = 2$, none beyond.

\textbf{9.} Each inner turn \omterm{def:b2:processor-and-pipelining:stall}{stalls} 2 at the \code{bge} and 2 at the \code{bltu}: $4 \times 120 = 480$. The stores are
not stalled, the \code{bge} having waited already. The printing loop's \code{bltu} follows its \code{addi}: $2
\times 16 = 32$. Each outer turn's \code{add t2, a0, t2} follows the \code{slli}: $2 \times 15 = 30$. Five pairs
occur once (the \code{sw ra} after \code{addi sp}; the two \code{la}, \code{addi} after \code{auipc}; the
\code{addi t1, t0, 64} after the second \code{la}; the first \code{blez} after \code{addi t0}): 10. Total $480 + 32 +
30 + 10 = 552$.

\textbf{10.} The taken branches each lose 2: 105 of the inner \code{bltu}, 55 of the swap test, 15 of the printing
loop and the last outer test, 176 in all, 352 cycles. The 16 \code{jal} (one \code{call}, 15 \code{j}) lose 1 each,
and \code{sort}'s \code{ret} loses 2; \code{main}'s \code{ret} is the last instruction and loses nothing. $352 + 16 +
2 = 370$.

\textbf{11.} $964 + 4 + 552 + 370 = 1890$ cycles, CPI $1890/964 \approx 1.961$.

\textbf{12.} The data \omterm{def:b2:processor-and-pipelining:hazard}{hazards}, 552 cycles against 370.

\textbf{13.} Only the load-use of \code{t4} by the \code{bge}, one cycle per inner turn: 120.

\textbf{14.} Move \code{addi t1, t1, 4} up between \code{lw t4} and the \code{bge}, and write the stores as
\code{sw t4, -4(t1)} and \code{sw t3, 0(t1)}: the \code{bge} is then two places after the load, and no \omterm{def:b2:processor-and-pipelining:stall}{stall}
remains, for the same number of instructions; 120 cycles saved. A further change saves loads too: after the compare
the larger value is at \code{p[1]}, so it can stay in a register as the next turn's \code{p[0]}, one load instead of
two per turn.

\textbf{15.} The swap test depends on the data and is wrong 41 times out of 120. The inner \code{bltu} is wrong at
each of the 15 exits and once at the very first turn, when its counter still says ``not taken'': 16, as
\cref{prop:b2:processor-and-pipelining:two-bit} predicts. The printing loop's is wrong at its first turn and at its
exit: 2. The outer \code{blez} is not taken 15 times and taken once, at the end: 1.

\textbf{16.} Its outcome is whether two neighbours are out of order, a property of the data, not of the loop's
structure; only as the sorted tail grows does it become regular.

\textbf{17.} $2 \times 60 + 16 + 2 = 138$.

\textbf{18.} $964 + 4 + 120 + 138 = 1226$ cycles, CPI $\approx 1.272$.

\textbf{19.} Single-cycle: $964 \times 800\,\text{ps} = 771.2$\,ns. Pipeline: $1890 \times 200 = 378.0$\,ns without
\omterm{def:b2:processor-and-pipelining:forwarding}{forwarding}, $1458 \times 200 = 291.6$\,ns with it, $1226 \times 200 = 245.2$\,ns with the predictor too.

\textbf{20.} $771.2 / 245.2 \approx 3.15$ times faster. The cycle is 4 times shorter but the CPI is 1.272 instead of 1:
$4 / 1.272 \approx 3.15$.

\textbf{21.} Pipelining acts on the cycle time (and raises the CPI from 1); \omterm{def:b2:processor-and-pipelining:forwarding}{forwarding} and prediction act on the CPI.
None changes the number of instructions.

\textbf{22.} Only the jumps would lose cycles, 18: $(964 + 4 + 120 + 18)/964 = 1106/964 \approx 1.147$.

\textbf{23.} With no \omterm{def:b2:processor-and-pipelining:stall}{stall} either: $986/964 \approx 1.023$, close to the ideal.

\textbf{24.} In the five-stage model, \omterm{def:b2:processor-and-pipelining:forwarding}{forwarding} and a two-bit predictor take the CPI of the \omterm{def:b2:processor-and-pipelining:stall}{bubble} sort \emph{from
1.961 to 1.272}: \omterm{def:b2:processor-and-pipelining:forwarding}{forwarding} removes 432 of the 552 \omterm{def:b2:processor-and-pipelining:stall}{stall} cycles, the predictor 232 of the 370 cycles lost to branches
and jumps, and the run takes 1226 cycles instead of 1890, 35\,\% fewer.
\end{solution}
